\documentclass[11pt,a4paper]{article}
\usepackage[T1]{fontenc}
\usepackage[utf8]{inputenc}
\usepackage[brazil]{babel}
\usepackage{lmodern}
\usepackage[margin=2.4cm]{geometry}
\usepackage{booktabs}
\usepackage{graphicx}
\usepackage{amsmath,amssymb}
\usepackage{caption}
\usepackage[colorlinks=false,hidelinks]{hyperref}
\usepackage{enumitem}
\usepackage{fancyhdr}

\graphicspath{{../../figures/}}
\captionsetup{font=small,labelfont=bf}
\setlist{itemsep=2pt,topsep=3pt}

\pagestyle{fancy}\fancyhf{}
\fancyhead[L]{\small UC day-ahead com VPL e VDC}
\fancyhead[R]{\small\thepage}
\renewcommand{\headrulewidth}{0.4pt}

\newcommand{\tabcustodia}{%
\begin{tabular}{@{}lrrrrrrrr@{}}
\toprule
dia & A (base) & B (+VPL) & C (+VDC) & D (ambos) & $\Delta_B$ & $\Delta_C$ & $\Delta_D$ & $I$ \\
\midrule
primavera (vento alto) & 684,7 & 611,3 & 660,2 & 586,2 & 73,4 & 24,5 & 98,5 & 0,6 \\
pico de verão & 2\,770,7 & 2\,698,7 & 2\,750,5 & 2\,679,8 & 72,1 & 20,3 & 90,9 & -1,5 \\
rampa solar & 1\,399,8 & 1\,313,1 & 1\,373,4 & 1\,291,5 & 86,7 & 26,4 & 108,3 & -4,8 \\
mediano & 1\,593,9 & 1\,519,8 & 1\,568,0 & 1\,493,3 & 74,1 & 25,9 & 100,6 & 0,6 \\
média & 1\,612,3 & 1\,535,7 & 1\,588,0 & 1\,512,7 & 76,6 & 24,3 & 99,6 & -1,3 \\
\bottomrule
\end{tabular}}

\newcommand{\taboperacional}{%
\begin{tabular}{@{}lrrrr@{}}
\toprule
métrica & A (base) & B (+VPL) & C (+VDC) & D (ambos) \\
\midrule
custo (k\$) & 1\,612,3 & 1\,535,7 & 1\,588,0 & 1\,512,7 \\
curtail. (MWh) & 20\,463,8 & 18\,907,2 & 19\,709,9 & 18\,056,9 \\
curtail. (\%) & 21,7 & 19,5 & 20,7 & 18,4 \\
partidas & 17,2 & 11,8 & 9,5 & 10,2 \\
h em Pmin & 260,8 & 249,2 & 246,5 & 242,5 \\
h congest. & 9,5 & 12,8 & 12,2 & 12,5 \\
spread (\$/MWh) & 4,1 & 5,0 & 5,1 & 5,4 \\
cap. virtual (MWh) & 0,0 & 269,4 & 0,0 & 282,9 \\
CO$_2$ (t) & 43\,206,4 & 42\,636,6 & 42\,595,3 & 42\,024,1 \\
\bottomrule
\end{tabular}}

\newcommand{\tabpenetracao}{%
\begin{tabular}{@{}lrrrrrrc@{}}
\toprule
VRE & A (k\$) & $\Delta_B$ & $\Delta_C$ & $\Delta_D$ & $I$ & ruído & $I<0$ \\
\midrule
30\% & 1\,948,7 & 56,8 & 18,2 & 66,7 & -8,3 & 2,6 & 4/4 \\
50\% & 1\,612,3 & 76,6 & 24,3 & 99,6 & -1,3 & 2,3 & 2/4 \\
70\% & 1\,444,7 & 81,8 & 29,0 & 107,7 & -3,1 & 2,3 & 4/4 \\
\bottomrule
\end{tabular}}

\newcommand{\tabsensibilidades}{%
\begin{tabular}{@{}lrrrr@{}}
\toprule
variação & $\Delta_B$ & $\Delta_C$ & $\Delta_D$ & $I$ \\
\midrule
base (200 MW, $f_d$=0,5) & 76,6 & 24,3 & 99,6 & -1,3 \\
VRE 30\% & 56,8 & 18,2 & 66,7 & -8,3 \\
VRE 70\% & 81,8 & 29,0 & 107,7 & -3,1 \\
VPL 100 MW & 38,0 & --- & 63,7 & 1,4 \\
VPL 400 MW & 141,5 & --- & 159,1 & -6,7 \\
VPL 650 MW & 185,8 & --- & 193,8 & -16,3 \\
$f_d$ = 0,2 & 78,6 & 11,3 & 88,6 & -1,3 \\
$f_d$ = 0,8 & 75,7 & 37,2 & 111,8 & -1,1 \\
janela $\times$0,5 & 76,0 & 8,0 & 82,7 & -1,3 \\
janela $\times$2 & 76,5 & 33,2 & 108,2 & -1,4 \\
migração $\times$0,5 & --- & 28,1 & 103,8 & -0,9 \\
migração $\times$2 & --- & 20,2 & 95,6 & -1,2 \\
reserva flexível & 76,8 & 27,0 & 103,6 & -0,1 \\
VPL explícito & 44,3 & --- & 69,7 & 1,2 \\
VPL dedicado & 40,0 & --- & 65,7 & 1,5 \\
\bottomrule
\end{tabular}}

\newcommand{\tabdecomposicao}{%
\begin{tabular}{@{}lrrrrr@{}}
\toprule
dia & $\Delta_B$ (par) & só montante & só jusante & pareamento & \% \\
\midrule
primavera (vento alto) & 73,4 & 33,0 & 37,3 & 3,2 & 4,4 \\
pico de verão & 72,1 & 43,1 & 28,5 & 0,5 & 0,7 \\
rampa solar & 86,7 & 40,7 & 44,4 & 1,5 & 1,8 \\
mediano & 74,1 & 33,7 & 33,4 & 6,9 & 9,3 \\
média & 76,6 & 37,6 & 35,9 & 3,1 & 4,1 \\
\bottomrule
\end{tabular}}

\newcommand{\tabsuperficie}{%
\begin{tabular}{@{}lrrr@{}}
\toprule
VPL/terminal & $f_d$=0,2 & $f_d$=0,5 & $f_d$=0,8 \\
\midrule
100 MW & -1,8 & 1,4 & -1,0 \\
200 MW & -1,3 & -1,3 & -1,1 \\
400 MW & -3,1 & -6,7 & -7,2 \\
\bottomrule
\end{tabular}}

\newcommand{\tabmensais}{%
\begin{tabular}{@{}lrrrrrrr@{}}
\toprule
dia & A (k\$) & $\Delta_B$ & $\Delta_C$ & $\Delta_D$ & $I$ & partidas A & partidas D \\
\midrule
mês 01 & 1\,126,5 & 78,9 & 20,4 & 98,9 & -0,4 & 16,0 & 16,0 \\
mês 02 & 1\,034,7 & 72,8 & 26,8 & 96,3 & -3,3 & 9,0 & 1,0 \\
mês 03 & 1\,226,5 & 72,4 & 33,7 & 107,9 & 1,7 & 16,0 & 13,0 \\
mês 04 & 1\,205,2 & 8,3 & 3,6 & 10,1 & -1,9 & 0,0 & 0,0 \\
mês 05 & 1\,570,3 & 70,1 & 13,4 & 86,3 & 2,9 & 12,0 & 15,0 \\
mês 06 & 1\,967,7 & 46,1 & 8,4 & 49,2 & -5,4 & 15,0 & 4,0 \\
mês 07 & 2\,756,6 & 70,5 & 16,9 & 84,0 & -3,5 & 31,0 & 15,0 \\
mês 08 & 2\,604,8 & 50,0 & 20,5 & 57,1 & -13,5 & 18,0 & 2,0 \\
mês 09 & 2\,223,2 & 33,4 & 15,1 & 41,9 & -6,6 & 13,0 & 2,0 \\
mês 10 & 1\,488,0 & 67,6 & 21,7 & 91,1 & 1,9 & 13,0 & 12,0 \\
mês 11 & 1\,243,3 & 80,8 & 21,4 & 101,6 & -0,6 & 16,0 & 11,0 \\
mês 12 & 1\,338,7 & 66,2 & 27,0 & 89,6 & -3,6 & 16,0 & 2,0 \\
média & 1\,648,8 & 59,8 & 19,1 & 76,2 & -2,7 & 14,6 & 7,8 \\
\bottomrule
\end{tabular}}

\newcommand{\tabcondicoes}{%
\begin{tabular}{@{}lrrrrrr@{}}
\toprule
dia & part. A & part. A herd. & part. D & part. D herd. & $I$ & $I$ herd. \\
\midrule
primavera (vento alto) & 7,0 & 21,0 & 4,0 & 7,0 & 0,6 & -5,0 \\
pico de verão & 26,0 & 41,0 & 14,0 & 19,0 & -1,5 & -9,7 \\
rampa solar & 21,0 & 25,0 & 7,0 & 11,0 & -4,8 & -10,6 \\
mediano & 15,0 & 17,0 & 16,0 & 7,0 & 0,6 & -0,6 \\
média & 17,2 & 26,0 & 10,2 & 11,0 & -1,3 & -6,5 \\
\bottomrule
\end{tabular}}

\newcommand{\tabbaseline}{%
\begin{tabular}{@{}lrrrrrr@{}}
\toprule
dia & $\Delta_C$ plano & $\Delta_C$ nativo & $\Delta_D$ plano & $\Delta_D$ nativo & $I$ plano & $I$ nativo \\
\midrule
primavera (vento alto) & 24,5 & 11,3 & 98,5 & 85,3 & 0,6 & -0,1 \\
pico de verão & 20,3 & 18,6 & 90,9 & 89,2 & -1,5 & -2,1 \\
rampa solar & 26,4 & 23,9 & 108,3 & 105,8 & -4,8 & -7,7 \\
mediano & 25,9 & 9,7 & 100,6 & 84,4 & 0,6 & 0,1 \\
média & 24,3 & 15,9 & 99,6 & 91,2 & -1,3 & -2,4 \\
\bottomrule
\end{tabular}}

\newcommand{\tabmigracao}{%
\begin{tabular}{@{}lrrrrrr@{}}
\toprule
dia & $\Delta_C$ base & $I$ base & $\Delta_C$ $\times$0,5 & $I$ $\times$0,5 & $\Delta_C$ $\times$2 & $I$ $\times$2 \\
\midrule
primavera (vento alto) & 24,5 & 0,6 & 30,3 & 0,6 & 20,7 & 1,4 \\
pico de verão & 20,3 & -1,5 & 23,9 & -1,7 & 14,1 & -2,2 \\
rampa solar & 26,4 & -4,8 & 31,2 & -3,4 & 21,5 & -4,5 \\
mediano & 25,9 & 0,6 & 27,2 & 0,8 & 24,5 & 0,6 \\
média & 24,3 & -1,3 & 28,1 & -0,9 & 20,2 & -1,2 \\
\bottomrule
\end{tabular}}

\newcommand{\tabsoc}{%
\begin{tabular}{@{}lrrrrrr@{}}
\toprule
dia & $\Delta_B$ livre & $\Delta_B$ 50\% & $\Delta_D$ livre & $\Delta_D$ 50\% & $I$ livre & $I$ 50\% \\
\midrule
primavera (vento alto) & 73,4 & 58,6 & 98,5 & 82,4 & 0,6 & -0,7 \\
pico de verão & 72,1 & 65,7 & 90,9 & 84,4 & -1,5 & -1,6 \\
rampa solar & 86,7 & 75,4 & 108,3 & 97,2 & -4,8 & -4,6 \\
mediano & 74,1 & 72,4 & 100,6 & 98,4 & 0,6 & 0,1 \\
média & 76,6 & 68,0 & 99,6 & 90,6 & -1,3 & -1,7 \\
\bottomrule
\end{tabular}}

\newcommand{\tableituras}{%
\begin{tabular}{@{}lrrrrrr@{}}
\toprule
dia & $\Delta_B$ emerg. & $I$ emerg. & $\Delta_B$ explíc. & $I$ explíc. & $\Delta_B$ dedic. & $I$ dedic. \\
\midrule
primavera (vento alto) & 73,4 & 0,6 & 48,5 & 1,7 & 34,0 & 3,2 \\
pico de verão & 72,1 & -1,5 & 33,0 & -2,1 & 31,5 & -2,2 \\
rampa solar & 86,7 & -4,8 & 53,9 & -2,1 & 53,2 & -1,8 \\
mediano & 74,1 & 0,6 & 41,8 & 7,2 & 41,2 & 6,8 \\
média & 76,6 & -1,3 & 44,3 & 1,2 & 40,0 & 1,5 \\
\bottomrule
\end{tabular}}

\newcommand{\tableiturasbig}{%
\begin{tabular}{@{}lrrrrrr@{}}
\toprule
dia & $\Delta_B$ emerg. & $I$ emerg. & $\Delta_B$ explíc. & $I$ explíc. & $\Delta_B$ dedic. & $I$ dedic. \\
\midrule
primavera (vento alto) & 138,4 & 0,1 & 93,4 & -1,9 & 68,8 & -0,6 \\
pico de verão & 121,2 & -16,5 & 59,2 & -3,6 & 54,9 & -4,1 \\
rampa solar & 165,0 & -10,7 & 97,9 & -6,8 & 96,7 & -6,6 \\
mediano & 141,5 & 0,2 & 88,7 & -1,9 & 85,9 & -0,8 \\
média & 141,5 & -6,7 & 84,8 & -3,6 & 76,5 & -3,0 \\
\bottomrule
\end{tabular}}

\newcommand{\tabreserva}{%
\begin{tabular}{@{}lrrrrrr@{}}
\toprule
dia & $\Delta_C$ base & $\Delta_D$ base & $I$ base & $\Delta_C$ $R$ comum & $\Delta_D$ $R$ comum & $I$ $R$ comum \\
\midrule
primavera (vento alto) & 24,5 & 98,5 & 0,6 & 24,1 & 98,3 & 0,7 \\
pico de verão & 20,3 & 90,9 & -1,5 & 19,7 & 90,2 & -1,6 \\
rampa solar & 26,4 & 108,3 & -4,8 & 26,4 & 107,9 & -5,1 \\
mediano & 25,9 & 100,6 & 0,6 & 25,6 & 100,2 & 0,5 \\
média & 24,3 & 99,6 & -1,3 & 23,9 & 99,1 & -1,4 \\
\bottomrule
\end{tabular}}

\newcommand{\numpareamentopct}{4}
\newcommand{\numparup}{37,6}
\newcommand{\numpardn}{35,9}
\newcommand{\numparpar}{76,6}
\newcommand{\nummensalcusto}{4,6}
\newcommand{\nummensalcurt}{15,4}
\newcommand{\nummensalsu}{46,9}
\newcommand{\nummensalIneg}{9}
\newcommand{\numredsubase}{41}
\newcommand{\numredsuhist}{58}
\newcommand{\numIbase}{-1,3}
\newcommand{\numIhist}{-6,5}
\newcommand{\numdCplano}{24,3}
\newcommand{\numdCnativo}{15,9}
\newcommand{\numsocdB}{68,0}
\newcommand{\numsocdD}{90,6}
\newcommand{\numsocI}{-1,7}
\newcommand{\numdBemerg}{76,6}
\newcommand{\numdBexplic}{44,3}
\newcommand{\numdBdedic}{40,0}
\newcommand{\numIdedic}{1,5}
\newcommand{\numdBbigemerg}{141,5}
\newcommand{\numdBbigexplic}{84,8}
\newcommand{\numdBbigdedic}{76,5}
\newcommand{\numIbigemerg}{-6,7}
\newcommand{\numIbigexplic}{-3,6}
\newcommand{\numIbigdedic}{-3,0}
\newcommand{\numcustoA}{1\,612,3}
\newcommand{\numpartA}{17,2}
\newcommand{\numcurtA}{21,7}
\newcommand{\numemisA}{43,2}
\newcommand{\numcustoB}{1\,535,7}
\newcommand{\numpartB}{11,8}
\newcommand{\numcurtB}{19,5}
\newcommand{\numemisB}{42,6}
\newcommand{\numcustoC}{1\,588,0}
\newcommand{\numpartC}{9,5}
\newcommand{\numcurtC}{20,7}
\newcommand{\numemisC}{42,6}
\newcommand{\numcustoD}{1\,512,7}
\newcommand{\numpartD}{10,2}
\newcommand{\numcurtD}{18,4}
\newcommand{\numemisD}{42,0}
\newcommand{\numredB}{4,7}
\newcommand{\numredC}{1,5}
\newcommand{\numredD}{6,2}
\newcommand{\numredcurtD}{11,8}
\newcommand{\numredsuD}{41}
\newcommand{\numncen}{320}
\newcommand{\numtmed}{10}
\newcommand{\numtmax}{1\,872}
\newcommand{\numgapmax}{0,120}
\newcommand{\numvdcespacial}{14,3}
\newcommand{\numvdctemporal}{50,7}
\newcommand{\numvdcit}{3\,181}

\title{\vspace{-1.5cm}\Large\textbf{Co-otimização de Virtual Power Lines e Virtual Data Centers\\
no Unit Commitment day-ahead de sistemas termoelétricos\\com alta penetração renovável}\\[10pt]
\large\normalfont Levando a Virtual Power Line da tese do Flávio para a operação day-ahead\\[3pt]
\small\itshape Co-Optimizing Virtual Power Lines and Virtual Data Centers in Day-Ahead Unit Commitment\\for Thermal-Dominated Systems with High Renewable Penetration}
\author{Guilherme Steilein \\ \small Doutorado --- PPGEE-UFPR}
\date{\small 11 de setembro de 2026}

\begin{document}
\maketitle
\thispagestyle{fancy}

\section*{Resumo}

A Virtual Power Line dimensionada no planejamento da tese do Flávio foi
levada para a operação, hora a hora, num Unit Commitment day-ahead com rede. No mesmo corredor
inseri um segundo recurso de flexibilidade que não existe na tese: data centers com carga adiável. O
objetivo é determinar se os dois recursos se somam ou se anulam parcialmente. Foram \numncen{} cenários
resolvidos com Gurobi. \textbf{O resultado é que eles competem}: juntos economizam de 1 a 11\,\% menos que a soma dos
benefícios isolados, embora façam coisas operacionalmente bem diferentes --- a bateria ataca o
curtailment, o data center ataca a ciclagem térmica.

% =====================================================================
\section{Motivação}

\subsection{A tese do Flávio}

A tese do Flávio trata de planejamento de expansão de geração e transmissão, de forma robusta e dinâmica,
e introduz o conceito central para este trabalho: a \textbf{linha de transmissão
virtual inter-área}. Em vez de reforçar fisicamente um tronco congestionado, coloca-se um par de baterias,
uma em cada terminal do corredor, e opera-se esse par de modo que ele emule capacidade de transferência
adicional. A bateria de montante absorve a energia que não conseguiu passar; a de jusante entrega essa
energia do outro lado quando é preciso. Do ponto de vista do sistema, o efeito é parecido com o de ter
mais corredor, só que instalável em meses em vez de quase uma década.

Quatro elementos da tese são reaproveitados aqui praticamente sem alteração:

\begin{itemize}
\item \textbf{O sistema teste estressado}: o RTS-GMLC com a carga da Área 2 aumentada em 40\,\% e uma
quantidade equivalente de geração renovável colocada na Área 3, nos nós candidatos que a tese identificou.
É essa modificação que transforma o tronco 318--223 no gargalo do sistema.
\item \textbf{O dimensionamento}: o modelo de planejamento da tese seleciona 650\,MW / 3\,800\,MWh para a
linha virtual. Uso esse valor como limite superior da minha varredura de tamanho.
\item \textbf{A regra de chaveamento}: a tese define o estado de cada bateria terminal em função do
sentido do fluxo no corredor e do estágio de uso da rede (alto ou baixo, pelos percentis da demanda
líquida). Implementei essa regra como variante e comparo com a operação livre.
\item \textbf{As métricas de rampa}: FRUS e FRDS, a sobra de rampa para cima e para baixo das unidades
ligadas, que a tese usa para argumentar sobre flexibilidade. Reporto as duas.
\end{itemize}

\subsection{Onde este trabalho se encaixa}

A tese do Flávio para no planejamento. Ela mostra que alocar armazenamento como linha virtual inter-área
adia o reforço do tronco por vários anos, e faz isso com uma camada de Unit Commitment acoplada de forma
fraca, só para verificar as rampas. A própria conclusão dele aponta que o passo seguinte é levar o
conceito para uma ferramenta de decisão operacional, onde o operador decide hora a hora.

\textbf{É exatamente esse o meu ponto de partida}: transformar a linha virtual de um ativo de planejamento
em um recurso de operação, dentro de um UC day-ahead completo, com rede, partidas, tempos mínimos e
rampas. Aí a pergunta deixa de ser ``quanto de bateria instalar'' e passa a ser ``o que essa bateria faz
em cada hora do dia, e quanto isso vale''.

\subsection{Contribuição deste trabalho}

No horizonte operacional, cabe identificar que outros agentes disputam a capacidade do mesmo corredor.
Daí a segunda peça, ausente na tese: os \textbf{data centers}. São hoje uma das cargas de maior
crescimento, e parte do que processam não tem restrição de horário nem de local --- treinamento de
modelos, análises em lote, backups. Essa parcela pode ser adiada algumas horas ou executada em outro
site.

Chamo esse conjunto de \textbf{VDC (Virtual Data Center)}. A analogia é direta: os dois recursos
aliviam o mesmo gargalo por caminhos opostos:

\begin{itemize}
\item a \textbf{VPL} move \emph{energia no tempo}, dentro do corredor: é uma solução pelo lado da oferta;
\item o \textbf{VDC} move \emph{carga no espaço e no tempo}: é uma solução pelo lado da demanda.
\end{itemize}

\textbf{A questão central:} usados em conjunto, esses recursos entregam mais ou menos que a soma dos
benefícios isolados? Não há na literatura consultada nenhum trabalho que os co-otimize no mesmo modelo. A métrica é direta:

\begin{equation}
I \;=\; \Delta_D - (\Delta_B + \Delta_C),
\qquad \Delta_X = \text{Custo}(A) - \text{Custo}(X),
\label{eq:I}
\end{equation}

onde A é o caso sem nada, B só com a VPL, C só com o VDC e D com os dois. Se $I>0$ eles se reforçam;
se $I<0$ um reduz o valor marginal do outro.

% =====================================================================
\section{Modelo e implementação}

\subsection{Ferramentas}

Não escrevi o Unit Commitment do zero. A base é o \textbf{PyPSA 1.3}, que monta a rede e gera
automaticamente as restrições clássicas do UC: balanço nodal, fluxo de potência linear, lógica de partida
e parada, tempos mínimos ligado e desligado, rampas e balanço de armazenamento. Por baixo dele está o
\textbf{linopy 0.9}, que constrói o problema algebricamente e conversa com o solver. A vantagem é que a
formulação padrão vem testada e auditável por terceiros; o que escrevi de fato são as duas camadas de
flexibilidade e três blocos de restrições adicionais.

\begin{table}[htbp]
\centering\small
\caption{O que cada peça faz.}
\begin{tabular}{@{}llp{8.4cm}@{}}
\toprule
Ferramenta & Versão & Papel \\
\midrule
PyPSA & 1.3.0 & rede, formulação do UC, fluxo DC, armazenamento, exportação NetCDF \\
linopy & 0.9.1 & montagem algébrica do MILP e interface com o solver \\
Gurobi & 13.0.3 & solver MILP da grade final (licença acadêmica da UFPR) \\
HiGHS & 1.15.1 & solver aberto, usado no desenvolvimento e na depuração \\
pandas / xarray & 3.0 / 2026.7 & manipulação de séries e de resultados \\
numpy / scipy & 2.5 / 1.18 & cálculo numérico e a DFT da seleção de dias \\
matplotlib & 3.11 & figuras \\
\bottomrule
\end{tabular}
\end{table}

Sobre o solver: a grade toda roda com \textbf{Gurobi 13.0.3 na licença acadêmica da UFPR}, ativada nesta
máquina em setembro. Não é detalhe menor. Como mostro na seção~\ref{sec:gap}, a interação que quero medir
é da mesma ordem de grandeza da tolerância que normalmente se aceita num MILP, e só foi possível resolver
a grade inteira a 0{,}05\,\% de gap por causa do Gurobi. Com o HiGHS, que é aberto e foi o que usei para
desenvolver, o mesmo conjunto levaria muito mais tempo e boa parte dos casos pararia num gap dez vezes
maior.

O código próprio tem cerca de \textbf{3\,700 linhas} em dez módulos, mais os testes, e tudo é encadeado
por um \texttt{Makefile}: uma linha de comando refaz desde a validação da rede até as figuras.

\subsection{Dados}

Nada foi inventado. São duas fontes, ambas clonadas de repositórios públicos:

\begin{itemize}
\item \textbf{RTS-GMLC}, do NREL e do consórcio GridMod (commit \texttt{3ece0d3}, de outubro de 2025).
Dele vêm a topologia (73 barras, 120 ramos), o parque gerador com curva de heat-rate por segmentos,
custos de partida, tempos mínimos e rampas de cada unidade, e as séries horárias de 2020 para carga,
eólica, solar fotovoltaica, solar de telhado e CSP. Uso sempre a versão \emph{day-ahead} das séries, que é
a que o operador teria em mãos na véspera.
\item \textbf{Phd\_Files}, o repositório do Flávio (commit \texttt{d1d593e}, de maio de 2024). Dele vêm os
nós candidatos de geração, que definem onde entra o acréscimo renovável da Área 3, além das planilhas de
custo que usei para conferência.
\end{itemize}

O que \emph{não} vem de dado público é o conjunto de data centers: tamanho, PUE e tarefas são premissas
minhas, geradas com semente fixa. Isso está declarado como limitação.

Alguns pontos da conversão do RTS-GMLC para o PyPSA que deram trabalho e vale registrar: os ramos vêm em
p.u. na base de 100\,MVA e precisam virar ohms; os ramos com relação de transformação viram
transformadores, não linhas, e isso muda o fluxo; a curva de heat-rate por segmentos é linearizada num
ponto de 80\,\% da potência nominal; o elo de corrente contínua vira um \emph{link}, enquanto as linhas CA
nunca podem virar \emph{link}, sob pena de perder a física do congestionamento. Validei tudo isso
comparando o fluxo do UC com um fluxo de potência linear estático: o erro de balanço nodal ficou na casa
de $10^{-12}$\,MW.

\subsection{Formulação}

O problema é um MILP. A função objetivo minimiza, ao longo das 24 horas, quatro parcelas:

\begin{equation}
\sum_{t}\Big[\underbrace{\sum_{g}\big(C^{\mathrm{E}}_g p_{g,t} + C^{\mathrm{SU}}_g v_{g,t}\big)}_{\text{combustível e partida}}
+ \underbrace{C^{\mathrm{curt}}\!\sum_{i}(\bar r_{i,t}-r_{i,t})}_{\text{curtailment}}
+ \underbrace{\sum_{b} C^{\mathrm{deg}}_b p^{\mathrm{dis}}_{b,t}}_{\text{degradação}}
+ \underbrace{\sum_{j,d} c_{j,d}\,\mathrm{PUE}_d\,x_{j,d,t}}_{\text{migração de carga}}\Big]
\end{equation}

As binárias são o estado, a partida e a parada de cada térmica ($u_{g,t}$, $v_{g,t}$, $w_{g,t}$):
\textbf{5\,256 binárias} por dia, o que faz deste um MILP de porte respeitável para 24 horas. As
contínuas são a geração, o fluxo em cada ramo, a carga e descarga das baterias e a potência de TI de cada
tarefa em cada site. No total, 14\,040 variáveis e 45\,910 restrições.

O PyPSA gera sozinho a lógica $u_{g,t}-u_{g,t-1}=v_{g,t}-w_{g,t}$, os tempos mínimos ligado e desligado,
os limites de potência, as rampas, o fluxo linear
$f_{\ell,t}=(\theta_{n_1}-\theta_{n_2})/x_\ell$ com $|f_{\ell,t}|\le F_\ell$, e o balanço em cada barra.

Escrevi à mão três blocos:

\begin{enumerate}
\item \textbf{Reserva girante}: a folga das unidades ligadas tem que cobrir 3\,\% da carga mais 5\,\% da
renovável disponível,
$\sum_g (u_{g,t}\overline{P}_g - p_{g,t}) \ge 0{,}03\,L_t + 0{,}05\sum_i \bar r_{i,t}$.
\item \textbf{Capacidade de cada data center}: a soma das tarefas rodando num site não pode passar da
capacidade flexível dele, $\sum_j \mathrm{PUE}_d\,x_{j,d,t} \le f_d C_d$.
\item \textbf{Rampa agregada de cada data center}: a carga de TI de um site não pode variar mais que
$\Delta_d$ de uma hora para a outra, $\left|\sum_j x_{j,d,t} - \sum_j x_{j,d,t-1}\right| \le \Delta_d$,
com $\Delta_d = 0{,}5\,f_d C_d/\mathrm{PUE}_d$. Esta é a restrição que estava declarada no artigo mas
aplicada por par tarefa--site na primeira versão do código; corrigi-la derrubou o valor do data center
de 26{,}4 para 24{,}3\,k\$/dia.
\item \textbf{A regra de chaveamento da tese}, em duas variantes: uma binária de sentido de fluxo no
corredor, com big-$M$, e outra que fixa o modo de cada bateria conforme o estágio de uso da rede. Na
variante \emph{explícita} as horas intermediárias ficam livres para arbitrar, o que corresponde ao
\emph{revenue stacking} da versão inter-área do capítulo 5 da tese. Na variante \emph{dedicada} a regra
vale nas 24 horas, que é a leitura estrita da seção 4.2.3. A comparação entre as duas está na
seção~\ref{sec:leituras}.
\end{enumerate}

Além dessas, o \emph{Link} do PyPSA traz uma rampa própria, que aplica o mesmo $\Delta_d$ a cada par
tarefa--site individualmente. Ela não é redundante em relação à rampa agregada: medindo as 168
soluções com VDC, fica ativa em todas elas e em 8{,}7\,\% dos pares (par tarefa--site, hora), um terço
deles fora da borda da janela. O efeito prático é que uma tarefa entra a meia potência na primeira
hora da sua janela. Por isso ela passou a constar da formulação do artigo como equação própria, e
não como detalhe de implementação.

A VPL em si são dois \emph{storage units} com estado de carga cíclico e eficiência $\sqrt{0{,}85}$ em cada
sentido, o que dá os 85\,\% de ida e volta da tese. Na formulação que uso como padrão \emph{não existe
nenhuma restrição ligando uma bateria à outra}: o comportamento de linha virtual emerge da posição delas e
dos preços. Volto a isso na seção~\ref{sec:decomp}, porque é um resultado por si só.

\textbf{Preços nodais.} Depois de resolver o MILP, fixo todas as binárias na solução e resolvo o problema
de novo como um LP. Os multiplicadores desse LP são os preços nodais. É o procedimento padrão de formação
de preço em mercado day-ahead: o preço sai com o \emph{commitment} já decidido.

\subsection{Modelo dos data centers}

\textbf{Quantos e onde.} São \numdcnsites{} sites, um por área, nas barras de maior carga de cada uma:
118 na Área 1, 218 na Área 2 e 318 na Área 3. A escolha da barra 318 não é neutra: é ela que põe o data
center no terminal de montante do corredor, disputando o mesmo gargalo que a VPL. Sem isso não haveria
interação para medir.

\begin{table}[htbp]
\centering\small
\caption{Composição dos três sites. A potência de TI é a elétrica dividida pelo PUE.}
\tabdcsites
\end{table}

\textbf{Potência e o papel do PUE.} Cada site tem \numdccap{}/\numdcnsites{}\,MW elétricos de capacidade,
\numdccap{}\,MW no total. O modelo separa potência de TI de potência elétrica: um megawatt de TI puxa
1{,}30\,MW da rede, e a diferença é refrigeração e perdas. Toda restrição de tarefa é escrita em MW de TI
e convertida para elétrico ao tocar a rede. Com $f_d = 0{,}5$, metade da capacidade é flexível, o que dá
\numdcflexit{}\,MW de TI por site; a outra metade entra como carga firme comum.

\textbf{Por que a capacidade flexível não roda cheia.} Existe um fator de utilização de \numdcutil{}, e
ele é essencial. Com utilização 1{,}0 a carga flexível teria de operar a plena potência as 24 horas e não
sobraria folga nenhuma para deslocar. \textbf{É a folga que cria a flexibilidade}, não a capacidade. Daí
a energia diária: \numvdcit{}\,MWh de TI adiáveis, ou \numdcelet{}\,GWh elétricos, mais
\numdcinflexdia{}\,GWh de carga inflexível, totalizando \numdctotaldia{}\,GWh, cerca de 10\,\% da carga do
sistema no dia de referência.

\textbf{As tarefas.} São 18 por dia, seis por site, de dois tipos que isolam os dois mecanismos de
flexibilidade.

\begin{table}[htbp]
\centering\small
\caption{As 18 tarefas por tipo. As \emph{batch} só rodam no site nativo e dão flexibilidade no tempo;
as migráveis rodam nos três sites e, com janela curta, dão flexibilidade no espaço.}
\tabdctarefas
\end{table}

Essa separação é deliberada. Modelos de UC que representam o data center como uma fração agregada da
energia diária não conseguem distinguir deslocamento no tempo de deslocamento no espaço, e são justamente
esses dois efeitos que se quer medir aqui.

\textbf{Como isso vira um modelo.} Cada tarefa se transforma em três componentes nativos do PyPSA, sem
nenhuma variável inventada:

\begin{enumerate}
\item uma \textbf{barra virtual}, que existe só para aquela tarefa;
\item um \textbf{\emph{link} de cada site elegível} até essa barra, com eficiência $1/\mathrm{PUE}$, custo
marginal igual ao custo de migração e janela de disponibilidade via \texttt{p\_max\_pu}: fora do intervalo
entre chegada e prazo, a potência máxima é zero. Escolher por qual \emph{link} a energia passa é escolher
\emph{onde} a tarefa roda;
\item um \textbf{\emph{store}} na barra virtual, cuja energia mínima salta para $E_j$ na hora do prazo.
Isso obriga a tarefa a terminar e permite que ela acumule energia em qualquer ritmo antes disso. Escolher
o ritmo é escolher \emph{quando} ela roda.
\end{enumerate}

O prazo passa a ser restrição dura sem nenhuma binária adicional, o que é a razão de o modelo resolver em
segundos. Quatro restrições governam o conjunto: o prazo, a capacidade do site descontada a carga
inflexível, a rampa agregada do site e a rampa de cada tarefa, as duas últimas de \numdcrampa{}\,MW de TI
por hora. O custo de migração vai de 5 a 20\,\$/MWh, é zero no site nativo e é assimétrico entre pares de
sites, o que impede a solução de migrar tudo de graça.

\textbf{O que o otimizador faz com isso.} No dia de vento alto, caso C, \numvdctemporal{}\,\% da energia
adiável muda de hora e \numvdcespacial{}\,\% muda de site. O atraso médio das tarefas é de
\numdcatraso{}\,h e o custo de migração pago no dia é de \numdcmigcusto{}\,k\$. Quando a bateria entra
junto, no caso D, o deslocamento temporal cai e o espacial sobe um pouco: é a competição pela folga do
corredor aparecendo no comportamento, não só no custo.

% =====================================================================
\section{Casos e cenários}

\subsection{Sistema teste}

RTS-GMLC, 73 barras, três áreas, 158 unidades, com a configuração estressada da tese. O corredor de
interesse é o tronco \textbf{318--223} (Área 3 para Área 2, 500\,MW). A penetração renovável é escalada
por um fator anual único, para 30, 50 ou 70\,\% da carga do ano.

\begin{figure}[htbp]
\centering
\includegraphics[width=0.6\linewidth]{fig1_map.pdf}
\caption{O sistema com o corredor em destaque, os dois terminais da VPL (triângulos) e os três data
centers (quadrados), um por área.}
\end{figure}

\subsection{Os quatro casos de estudo}

\begin{itemize}
\item \textbf{A} --- base, sem nenhum recurso flexível;
\item \textbf{B} --- A mais a VPL (200\,MW / 4\,h por terminal, 85\,\% de eficiência ida-e-volta, como na
tese);
\item \textbf{C} --- A mais o VDC (três sites de 200\,MW, PUE 1{,}30, metade da capacidade flexível);
\item \textbf{D} --- A com os dois.
\end{itemize}

\textbf{Equalização de energia entre os casos.} Nos casos A e B os data centers permanecem no sistema,
como carga inflexível, consumindo exatamente a mesma energia diária que consomem em C e D. Sem essa
equalização, a comparação mediria o valor de \emph{retirar carga} do sistema em vez do valor da
\emph{flexibilidade} dela. O efeito sobre o resultado é da ordem de 40\,\% do valor atribuído ao data
center.

\subsection{Grade de cenários}

Foram \numncen{} resoluções, todas com Gurobi a 0{,}05\,\% de gap (mediana de \numtmed{}\,s por caso):

\begin{itemize}
\item quatro dias representativos $\times$ quatro casos $\times$ doze variações;
\item uma varredura cruzada de tamanho da VPL contra fração flexível do data center;
\item cada bateria terminal rodando sozinha, para separar o que é ``par'' do que é ``bateria'';
\item doze dias mensais, um por mês, para ver se os resultados valem fora dos dias extremos;
\item o dia anterior resolvido, para herdar as condições iniciais das térmicas;
\item um baseline alternativo para os data centers inflexíveis.
\end{itemize}

Os quatro dias vieram de uma seleção automática sobre a série anual, usando a mesma caracterização de
rampa da demanda líquida que a tese usa: o dia de maior excedente eólico (12/03), o pico anual de carga
(26/08), a maior rampa de três horas (14/01) e um dia útil de energia mediana (16/03).

% =====================================================================
\section{Resultados}

\subsection{Valor de cada recurso}

\begin{table}[htbp]
\centering\small
\caption{Custo operacional (k\$/dia) e benefícios, com 50\,\% de renovável. $\Delta$ é a economia em
relação ao caso A; $I$ é a interação da equação~\eqref{eq:I}.}
\tabcustodia
\end{table}

Na média dos quatro dias, a VPL sozinha economiza \numparpar{}\,k\$/dia (\numredB{}\,\% do custo), o
data center flexível sozinho \numdCplano{}\,k\$ (\numredC{}\,\%) e os dois juntos \numredD{}\,\%. A
interação média é \numIbase{}\,k\$/dia, ligeiramente negativa. A interação varia de sinal entre os dias, o que indica desde já a necessidade de tratamento numérico
cuidadoso.

\subsection{Serviços operacionais distintos}

\begin{table}[htbp]
\centering\small
\caption{Métricas operacionais, médias dos quatro dias.}
\taboperacional
\end{table}

Os dois recursos não são intercambiáveis do ponto de vista operacional:

\begin{itemize}
\item \textbf{A VPL ataca o curtailment}: derruba de \numcurtA{}\,\% para \numcurtB{}\,\%. Nas partidas ela
mexe pouco.
\item \textbf{O VDC ataca a ciclagem}: as partidas caem de \numpartA{} para \numpartC{} por dia, e as horas
em mínimo técnico também. No curtailment ele quase não mexe.
\item Juntos chegam a \numpartD{} partidas e ao menor curtailment.
\item As \emph{horas congestionadas aumentam} com qualquer um dos recursos, de 9{,}5 para 10{,}8--12{,}8.
O comportamento é esperado: os recursos passam a usar o corredor em mais horas para mover energia ou
carga através dele. O corredor fica mais utilizado, não menos.
\end{itemize}

\begin{figure}[htbp]
\centering
\includegraphics[width=0.88\linewidth]{fig2_dispatch.pdf}
\caption{Despacho por combustível nos casos A e D, no dia de vento alto. O curtailment é a área
hachurada; no caso D ela diminui e o perfil térmico fica mais suave.}
\end{figure}

\subsection{Operação hora a hora}

\begin{figure}[htbp]
\centering
\begin{minipage}{0.48\linewidth}\centering
\includegraphics[width=\linewidth]{fig3_corridor_vpl.pdf}
\end{minipage}\hfill
\begin{minipage}{0.48\linewidth}\centering
\includegraphics[width=\linewidth]{fig4_vdc_heatmap.pdf}
\end{minipage}
\caption{À esquerda, o fluxo no corredor e o estado de carga das duas baterias, na variante com a regra de
chaveamento da tese, que deixa visíveis os estágios de alto e baixo uso. À direita, a carga dos data
centers por site e hora, nos casos C e D.}
\end{figure}

A bateria de montante carrega de madrugada e de novo entre 12\,h e 16\,h, enquanto o corredor está
saturado e há vento sendo cortado na Área 3. A de jusante carrega pelo corredor no vale do meio-dia e
descarrega no pico da noite, entre 18\,h e 20\,h, quando o fluxo cai de 500 para cerca de 270\,MW. É o
comportamento descrito na tese, e ele aparece mesmo sem impor a regra: na operação livre o otimizador
converge para o mesmo padrão.

Do lado dos data centers, a carga adiável se concentra entre 6\,h e 17\,h e o pico da noite fica quase
vazio nos três sites.

\textbf{Por que eles competem, se fazem coisas diferentes?} Porque precisam do mesmo insumo escasso, que
é a \emph{folga do corredor nas horas em que ele não está saturado}. A bateria de jusante só pode ser
carregada com energia que atravessa o corredor. A tarefa migrada para a Área 3 só tem valor porque consome
ali a energia que seria cortada. Quando um dos dois ocupa essa folga, sobra menos para o outro. Nos dados
isso aparece com clareza: no caso D a capacidade virtual da VPL e o deslocamento temporal do data center
ficam ambos menores do que nos casos isolados.

Um resultado contrário à hipótese inicial: o \emph{spread} de preço entre as pontas do corredor
\textbf{não colapsa}. Ele permanece em 21--23\,\$/MWh nas horas saturadas nos quatro casos, apenas se
desloca de hora. A competição, portanto, é por capacidade física do corredor, e não pela eliminação de um
sinal de preço.

\subsection{Interação por nível de penetração}

\begin{table}[htbp]
\centering\small
\caption{Valor marginal e interação por nível de penetração, médias dos quatro dias (k\$/dia). A coluna
``ruído'' é a soma dos gaps absolutos dos quatro casos, o limite abaixo do qual $I$ não se distingue da
tolerância do solver.}
\tabpenetracao
\end{table}

\begin{itemize}
\item A interação é negativa nos três níveis.
\item \textbf{A competição é mais forte a 30\,\%}: o corredor satura em poucas horas, quatro em média, e
os dois recursos disputam a mesma janela curta. O benefício conjunto fica 11{,}1\,\% abaixo da soma.
\item A 50 e 70\,\% o corredor satura em 9 a 11 horas, há mais espaço para atuarem em horas diferentes, e
a interação encolhe, sem mudar de sinal.
\item A hipótese inicial era a oposta: com muita renovável, a bateria ficaria limitada por energia e o
data center por prazo, tornando-os complementares. \textbf{Os dados não sustentam essa hipótese.}
\end{itemize}

Olhando as 52 combinações de dia e variação da grade principal em que os quatro casos foram
resolvidos, $I$ é negativo em 34 e supera duas vezes o ruído do gap em 19 delas, sendo 15 de
competição e 4 de complementaridade.

\begin{figure}[htbp]
\centering
\includegraphics[width=0.72\linewidth]{fig7_complementarity.pdf}
\caption{Interação por dia e penetração. Barra negativa indica competição entre os recursos.}
\end{figure}

\subsection{Sensibilidades}

\begin{table}[htbp]
\centering\small
\caption{Efeito de cada variação sobre os benefícios e sobre a interação (k\$/dia, médias dos quatro
dias). Onde aparece ``---'', o eixo não se aplica àquele caso.}
\tabsensibilidades
\end{table}

\begin{itemize}
\item \textbf{Tamanho da VPL}: o valor cresce, mas com retorno marginal decrescente --- 38, 32 e 18\,k\$
por cada 100\,MW adicionais, ao passar de 100 para 200, 400 e 650\,MW. Portanto \textbf{os 650\,MW selecionados pelo
planejamento da tese ainda agregam valor na operação}, a cerca de metade da taxa marginal inicial. É um
ponto de consistência direta entre os dois trabalhos.
\item \textbf{Fração flexível do data center}: com 0{,}2 ele quase não vale nada, porque paga custo de
migração para deslocar pouca energia; com 0{,}8 vale 37\,k\$.
\item \textbf{Largura da janela}: dobrar a janela das tarefas \emph{batch} leva o valor de 24 para 33\,k\$.
O prazo é restrição de primeira ordem, tanto quanto a quantidade de carga flexível.
\item \textbf{A regra de chaveamento da tese}: na variante explícita custa 2{,}1\,\% do custo
operacional, ou seja, abre mão de 32 dos 77\,k\$ de benefício da VPL, porque proíbe a bateria de arbitrar
contra o sentido do fluxo nas horas estagiadas. Na variante dedicada, que fecha também as horas
intermediárias, a perda sobe para 37 dos 77\,k\$. A seção~\ref{sec:leituras} detalha.
\item \textbf{Reserva provida pelos recursos flexíveis}: efeito quase nulo, entre 0{,}01 e 0{,}2\,\%. A
reserva raramente é restritiva depois que o data center já achatou a demanda líquida.
\end{itemize}

\begin{figure}[htbp]
\centering
\includegraphics[width=\linewidth]{fig9_sizing.pdf}
\caption{Custo, curtailment e partidas em função do tamanho da VPL, da fração flexível e da largura da
janela. As linhas pontilhadas são os casos aos quais o eixo não se aplica.}
\end{figure}

\subsection{Superfície de interação: tamanho da VPL $\times$ flexibilidade}

Como a competição depende de quanto cada recurso consegue absorver da oportunidade do corredor, cruzei os
dois eixos de dimensionamento.

\begin{table}[htbp]
\centering\small
\caption{Interação $I$ (k\$/dia) em função do tamanho da VPL e da fração flexível, médias dos quatro dias.
O ruído do gap é de cerca de 2{,}5\,k\$ em todas as células.}
\tabsuperficie
\end{table}

\begin{figure}[htbp]
\centering
\includegraphics[width=\linewidth]{fig8_cross_vpl_fd.pdf}
\caption{Superfície de interação, por dia e na média. Vermelho é competição.}
\end{figure}

\begin{itemize}
\item Com VPL de 100 a 200\,MW a interação média fica \emph{dentro do ruído} em qualquer fração flexível.
Ou seja, nesse porte, dá para tratar os dois recursos como aditivos.
\item \textbf{Com 400\,MW ela fica claramente negativa}: $-3{,}1$, $-6{,}7$ e $-7{,}2$\,k\$/dia, o que
equivale a 19 a 28\,\% do próprio benefício do data center. Com os 650\,MW da tese chega a
$-16{,}3$\,k\$/dia.
\item Complementaridade significativa aparece numa célula só, e num dia só: VPL de 100\,MW com metade da carga
flexível, no dia mediano, com $+7{,}0$\,k\$ (2,4 vezes o ruído). Faz sentido: a bateria é pequena demais
para esgotar a oportunidade do corredor, e o data center preenche as horas que sobram.
\end{itemize}

\subsection{Contribuição do pareamento dos terminais}
\label{sec:decomp}

Esta verificação toca diretamente o conceito da tese. Na formulação adotada como padrão, as duas
baterias não têm nenhuma restrição ligando uma à outra: o comportamento de linha virtual \emph{emerge} da
localização e dos preços. Cabe então separar o que vem do pareamento e o que vem simplesmente de haver
duas baterias bem posicionadas. Para isso, cada terminal foi resolvido isoladamente.

\begin{table}[htbp]
\centering\small
\caption{Decomposição do benefício da VPL (k\$/dia). ``Pareamento'' é o benefício do par menos a soma dos
benefícios de cada bateria sozinha.}
\tabdecomposicao
\end{table}

Cada bateria sozinha vale \numparup{} (montante) e \numpardn{}\,k\$ (jusante); o par vale
\numparpar{}\,k\$. \textbf{O pareamento acrescenta só \numpareamentopct{}\,\%.} Ou seja, praticamente todo
o benefício vem de ter armazenamento \emph{nos terminais do corredor certo}, coordenado ou não.

O resultado não enfraquece a proposta da tese: reforça a sua componente de planejamento. O que a tese
resolve é \emph{onde} alocar o armazenamento, e é essa decisão de localização que carrega quase todo o
valor. O que estes números indicam é que, uma vez bem localizadas, as baterias entregam o benefício sem
necessidade de coordenação explícita --- ao menos num modelo determinístico e sem contingências, que é o
caso aqui.

O mesmo vale para o outro resultado: a regra de chaveamento da tese \emph{custa} de 2 a 2{,}4\,\% do
custo operacional em vez de agregar, conforme se liberem ou não as horas intermediárias. A leitura precisa de cautela. A regra é um protocolo simples, aplicável sem rodar um UC completo,
e o que se mede aqui é o custo de oportunidade dela em ambiente idealizado. Num modelo com contingências,
obrigações de transmissão ou incerteza, uma regra que garante o comportamento de linha virtual pode valer
mais do que a arbitragem livre permitida ao otimizador.

\subsection{A VPL na definição da tese}
\label{sec:leituras}

Vale separar a VPL enquanto ativo da VPL enquanto regra, porque a tese define pela regra. Na seção 4.2.3
cada ESS opera permanentemente em estado de carga ou de descarga, e esse estado é determinado por dois
indicadores: o estágio de demanda líquida e o sentido do fluxo no corredor. A Tabela 4.3 da tese fecha as
quatro combinações. A seção 4.2.13 diz o mesmo de forma direta: ao contrário de um elemento físico de
transmissão, governado pelas equações de fluxo, a VTL exige restrições lógicas que acoplem estado do
armazenamento, sentido do fluxo e nível de utilização da rede. E na taxonomia de serviços do capítulo 4 o
armazenamento alocado a uma VTL emula transferências \emph{exclusivamente}, sem prestar nenhum outro
serviço; arbitragem é uma categoria separada, com regra própria. O empilhamento de receitas só aparece na
versão inter-área do capítulo 5, e mesmo lá a função de transferência fica reservada a estágios
específicos.

Isso expõe uma folga na variante que eu vinha chamando de explícita: ela aplica a regra nos estágios
extremos e libera as horas intermediárias para arbitragem, o que corresponde ao tratamento do capítulo 5 e
não ao do capítulo 4. Implementei então uma terceira variante, \textbf{dedicada}, em que a regra vale em
todas as horas e não sobra nenhuma hora livre.

\begin{table}[htbp]
\centering\small
\caption{As três leituras da VPL: benefício e interação (k\$/dia, médias dos quatro dias).}
\tableituras
\end{table}

O protocolo estrito custa caro: a VPL cai de \numdBemerg{} para \numdBdedic{}\,k\$/dia, ou seja,
praticamente metade do benefício. Mas quase todo esse custo já aparece na variante explícita, que entrega
\numdBexplic{}\,k\$. \textbf{Fechar as horas intermediárias custa só mais 4\,k\$}; o grosso da conta vem da
própria regra de estágio. A capacidade virtual acompanha, caindo de 269 para 191 e depois 174\,MWh, e o
\emph{throughput} de 1\,598 para 1\,360 e 1\,129\,MWh.

No dia de vento alto, sob a regra estrita, \textbf{a bateria de montante fica parada as 24
horas}. Todo o benefício do dia vem da bateria de jusante. Nos outros três dias ela opera de 5 a 12 horas,
contra 11 a 18 na formulação livre.

A razão é direta. No corredor estudado o fluxo nunca inverte: é sempre de montante para jusante, nas 96
horas. Com o sentido fixo, a Tabela 4.3 deixa a bateria de montante com apenas duas opções, carregar nas
horas de alto uso e descarregar nas de baixo uso. É o inverso do que a economia pede, porque o excedente
renovável de montante ocorre ao meio-dia, que é justamente hora de baixa demanda líquida. No dia de vento
alto, em que esse descasamento é mais forte, o otimizador simplesmente não opera essa bateria.

Sobre a interação, a leitura precisa de cuidado. No porte base ela vai de \numIbase{} na formulação livre
para \numIdedic{}\,k\$/dia na dedicada, trocando de sinal. \textbf{Mas as três ficam dentro do ruído do
gap}, que é de 2,3 a 2,5\,k\$ nas três variantes. Ou seja, no porte base nenhuma das leituras produz
interação mensurável, o que é consistente com o que já se sabia: o sinal de competição aparece com VPL
grande ou com penetração baixa, não no ponto base.

Por isso repeti a comparação a 400\,MW por terminal, que é onde a competição é mensurável na formulação
livre. Esta é a verificação que de fato interessa, porque testa se a conclusão do trabalho depende de eu
ter deixado as baterias livres.

\begin{table}[htbp]
\centering\small
\caption{As três leituras com VPL de 400\,MW por terminal (k\$/dia, médias dos quatro dias).}
\tableiturasbig
\end{table}

\textbf{A competição sobrevive ao protocolo da tese.} Com 400\,MW a interação é
\numIbigemerg{}\,k\$/dia na formulação livre, \numIbigexplic{} na explícita e \numIbigdedic{} na
dedicada. Nas duas variantes com regra ela é \textbf{negativa nos quatro dias}, contra dois dos quatro
na formulação livre.
O padrão é coerente com o mecanismo. Amarrar a bateria reduz o quanto ela captura da oportunidade do
corredor, e por isso ela também disputa menos: o benefício cai de \numdBbigemerg{} para
\numdBbigdedic{}\,k\$/dia e a competição cai junto, de \numIbigemerg{} para \numIbigdedic{}, deixando de
superar duas vezes o ruído na média, embora o sinal fique mais consistente dia a dia. Em proporção ao
benefício do data center, a competição vai de 28\,\% na formulação livre para 12\,\% na dedicada.
\textbf{O sentido do resultado não depende de qual leitura se adote; a magnitude, sim.}


Três observações sobre por que nada disso enfraquece a tese:

\begin{enumerate}
\item \textbf{O estágio da tese é anual, o meu é diário.} Lá a discretização da demanda líquida vem da
curva de duração anual dentro de um modelo de planejamento. Aplicá-la hora a hora num único dia é um uso
mais severo do que a regra foi desenhada para suportar.
\item \textbf{A regra pressupõe fluxo que inverte.} A lógica da Tabela 4.3 só usa os dois terminais quando
o sentido do fluxo muda ao longo do dia. Num corredor unidirecional saturado, o terminal de montante perde
função sob a regra.
\item \textbf{A métrica da tese é outra.} Lá a VTL é avaliada pelo investimento em linha que ela adia, não
pelo custo operacional de um dia. Estes números não medem aquilo.
\end{enumerate}

O que levo daqui é que os três valores formam um intervalo, e não um veredito: a VPL vale
\numdBdedic{}\,k\$/dia sob o protocolo estrito, \numdBexplic{} com as horas intermediárias liberadas e
\numdBemerg{} se as baterias forem totalmente livres. O primeiro número é o que corresponde à definição da
tese; o último é um limite superior do que armazenamento de corredor pode entregar.

\subsection{Validade fora dos dias extremos}

\begin{table}[htbp]
\centering\small
\caption{Um dia útil de energia mediana por mês, configuração base (k\$/dia).}
\tabmensais
\end{table}

Nos doze dias mensais o uso conjunto reduz o custo em \nummensalcusto{}\,\%, o curtailment em
\nummensalcurt{}\,\% e as partidas em \nummensalsu{}\,\%, contra 
umredD{}\,\%, 
umredcurtD{}\,\% e 
umredsuD{}\,\% nos quatro dias
representativos. \textbf{A interação é negativa em \nummensalIneg{} dos 12 dias}, com média de
$-2{,}7$\,k\$/dia. Os quatro dias escolhidos são, portanto, otimistas em magnitude, como esperado por serem dias extremos,
mas o sinal da competição se mantém ao longo do ano.

% =====================================================================
\section{Verificações de robustez}

Como a conclusão é contraintuitiva, sete hipóteses de modelagem que poderiam produzi-la artificialmente
foram testadas isoladamente.

\subsection{Precisão do solver}
\label{sec:gap}

Esta foi a questão que mais mudou a forma de rodar o trabalho, então vale detalhar.

\textbf{O que é o \emph{gap}.} Num MILP o solver mantém duas referências: a melhor solução inteira já
encontrada, que é um limite superior do custo, e o melhor limite inferior que a relaxação permite provar.
A distância relativa entre as duas é o \emph{gap}. O solver para quando essa distância cai abaixo da
tolerância pedida. Ou seja, com tolerância de 0{,}5\,\% eu não sei onde está o ótimo: sei apenas que o
custo que reportei está em algum ponto até 0{,}5\,\% acima dele.

\textbf{Por que isso é crítico aqui.} A interação da equação~\eqref{eq:I} pode ser reescrita como

\begin{equation}
I \;=\; \Delta_D - (\Delta_B + \Delta_C) \;=\; (B + C) - (A + D),
\end{equation}

isto é, uma combinação de quatro custos com sinais opostos. O erro de truncamento de cada um é sempre
positivo (o custo reportado nunca é menor que o ótimo), mas como dois entram somando e dois subtraindo,
\textbf{os erros não se cancelam: no pior caso eles se acumulam}.

\begin{table}[htbp]
\centering\small
\caption{Erro que a tolerância do solver admite, para um custo típico de 1{,}56\,M\$/dia por caso.}
\begin{tabular}{@{}lrrl@{}}
\toprule
Tolerância & Erro por caso & Erro possível em $I$ & Comparação \\
\midrule
0{,}5\,\% (usual) & até 7{,}8\,k\$ & $\pm$15{,}6\,k\$ & maior que qualquer interação medida \\
0{,}05\,\% (adotada) & até 0{,}8\,k\$ & $\pm$1{,}6\,k\$ & menor que a maioria das interações \\
\bottomrule
\end{tabular}
\end{table}

As interações que apareceram na grade vão de menos de 1 até 20\,k\$/dia. Com tolerância de 0{,}5\,\%, portanto,
qualquer valor medido estaria dentro da margem de erro do próprio solver --- eu não teria como afirmar
sequer o \emph{sinal} da interação.

\textbf{Um exemplo concreto.} Resolvendo o caso D do dia de vento alto duas vezes, só mudando a
tolerância:

\begin{itemize}
\item com alvo de 0{,}5\,\%, o Gurobi para em 7{,}3\,s e reporta 585\,526\,\$;
\item com alvo de 0{,}05\,\%, para em 7{,}7\,s e reporta 583\,185\,\$.
\end{itemize}

São 2\,341\,\$ de diferença no mesmo caso, pela mesma máquina, só por causa da tolerância. A interação
medida naquele dia é de 1\,216\,\$. \textbf{O erro de um único dos quatro casos seria o dobro do efeito
que eu quero medir.}

\textbf{E por que o Gurobi?} Aqui é preciso ser preciso: o gap é uma escolha de quem roda, não uma
propriedade do solver. O HiGHS também aceita 0{,}05\,\%. A diferença está no custo computacional de
chegar lá, porque o esforço para fechar os últimos décimos de percentual cresce de forma muito não
linear: os primeiros pontos percentuais vêm de heurísticas rápidas, e o resto exige explorar a árvore de
\emph{branch-and-bound}. Medi isso nos mesmos dois casos:

\begin{table}[htbp]
\centering\small
\caption{Tempo para fechar cada tolerância, mesmo caso e mesma máquina (8 threads).}
\begin{tabular}{@{}llrrr@{}}
\toprule
Caso & Solver & alvo 0{,}5\,\% & alvo 0{,}05\,\% & fator \\
\midrule
A (base)  & HiGHS  & 9{,}1\,s  & 80{,}5\,s & 8{,}8$\times$ \\
A (base)  & Gurobi & 5{,}4\,s  & 9{,}3\,s  & 1{,}7$\times$ \\
D (ambos) & HiGHS  & 15{,}8\,s & 28{,}5\,s & 1{,}8$\times$ \\
D (ambos) & Gurobi & 7{,}3\,s  & 7{,}7\,s  & 1{,}1$\times$ \\
\bottomrule
\end{tabular}
\end{table}

Nesses dois casos o HiGHS até fecha o 0{,}05\,\%, mas leva de 2 a 9 vezes mais tempo, e chega a ser
9 vezes mais lento que o Gurobi na mesma tolerância. O problema aparece na escala da grade: são
\numncen{} cenários, e os casos difíceis (os de 70\,\% de renovável, em que o Gurobi já chega a bater o
limite de 30 minutos) simplesmente não fechariam. Como referência, na grade antiga que rodei inteira em
HiGHS com tolerância de 0{,}5\,\%, 13 dos 116 casos pararam \emph{no teto} da tolerância sem provar
otimalidade --- ou seja, já estavam no limite do que aquele solver entregava.

\textbf{Resumo.} Com o HiGHS a 0{,}5\,\% o sinal de $I$ era inconclusivo em \emph{todas} as 48
combinações da grade de então, e chegou a inverter entre as duas tolerâncias em cinco delas. A licença acadêmica da UFPR
foi o que tornou viável rodar a grade inteira a 0{,}05\,\% em algumas horas. A lição que levo daqui é
metodológica: estudos que medem ``sinergia'' entre tecnologias como diferença de objetivos de MILP
precisam reportar o gap, senão podem estar medindo tolerância de solver. Passei a reportar o ruído ao
lado de cada interação.

\subsection{Condições iniciais das térmicas}

O RTS-GMLC não traz histórico de operação, então é preciso decidir quais unidades começam o dia ligadas.
Como partidas são métrica central, a hipótese foi testada: o dia anterior a cada dia representativo foi
resolvido e o estado final das térmicas herdado.

\begin{table}[htbp]
\centering\small
\caption{Regra padrão contra estado herdado do dia anterior.}
\tabcondicoes
\end{table}

As partidas absolutas sobem bastante, de 17{,}2 para 26{,}0 no caso A, porque no estado real há menos
unidades ligadas à meia-noite. Mas \textbf{a redução obtida pelo uso conjunto sobe também, de
\numredsubase{}\,\% para \numredsuhist{}\,\%}, e a competição fica \emph{mais} forte: a interação vai de
\numIbase{} para \numIhist{}\,k\$/dia, negativa nos quatro dias. A configuração adotada como padrão é,
portanto, a conservadora.

\subsection{Baseline dos data centers nos casos A e B}

A carga plana usada nos casos sem flexibilidade é mais suave que um perfil realista de tarefas chegando
ao longo do dia, o que poderia inflar o valor atribuído à flexibilidade. Foi testado um baseline mais restritivo: as
mesmas tarefas executadas no site nativo assim que chegam, por ordem de prazo.

\begin{table}[htbp]
\centering\small
\caption{Baseline plano contra baseline nativo (k\$/dia).}
\tabbaseline
\end{table}

O valor do data center cai de \numdCplano{} para \numdCnativo{}\,k\$, um terço menor. A interação muda de
magnitude, de $-1{,}3$ para $-2{,}4$\,k\$, mas \textbf{não muda de sinal em nenhum dos quatro dias}. O baseline plano é mantido como principal e este entra como sensibilidade.

\subsection{Custo de migração}

Como a migração entre sites é parte central do conceito de VDC, foram testados metade e o dobro do
custo.

\begin{table}[htbp]
\centering\small
\caption{Sensibilidade ao custo de migração (k\$/dia).}
\tabmigracao
\end{table}

O valor do data center varia $+15$ e $-16$\,\%, mas a interação fica em $-0{,}9$ e $-1{,}2$\,k\$ contra
$-1{,}3$ na base. \textbf{A conclusão não depende desse parâmetro}, que é o de maior arbitrariedade no
modelo.

\subsection{Estado de carga inicial da VPL}

As baterias operam com estado de carga cíclico: o nível ao fim das 24\,h tem que igualar o do início,
mas o otimizador escolhe qual é esse nível. Isso dá ao modelo uma liberdade que um operador real não
tem, já que o estado à meia-noite vem do dia anterior. Para medir o peso dessa hipótese, a grade foi
repetida com o estado de carga fixado em 50\,\% no início e no fim.

\begin{table}[htbp]
\centering\small
\caption{Estado de carga cíclico livre contra ancorado em 50\,\% (k\$/dia).}
\tabsoc
\end{table}

O benefício da VPL cai de \numparpar{} para \numsocdB{}\,k\$, 11\,\% menos, e o do uso conjunto de
\numcustoA{} --- na verdade de 99{,}6 para \numsocdD{}\,k\$. A perda vem toda do dia de vento alto, onde
a liberdade de começar o dia com a bateria cheia valia 15\,k\$. \textbf{A interação não muda de sinal}:
vai de \numIbase{} para \numsocI{}\,k\$/dia, ou seja, ancorar o estado inicial torna a competição
ligeiramente mais forte, não mais fraca. A hipótese cíclica livre é, portanto, otimista quanto ao valor
absoluto da VPL e conservadora quanto à conclusão do trabalho.

\subsection{O requisito de reserva era o mesmo nos quatro casos?}

Esta verificação nasceu de uma observação externa e valeu a pena. O requisito de reserva girante é
$R_t = 0{,}03\,L_t + 0{,}05\sum_i \bar r_{i,t}$, e $L_t$ é a soma das cargas do modelo. Nos casos A e B o
data center entra como carga plana e, portanto, conta inteiro em $L_t$. Nos casos C e D a parte flexível
passa por \emph{links} e não é uma carga: só a parte inflexível conta. Ou seja, \textbf{C e D tinham um
requisito de reserva menor que A e B}, e parte do benefício atribuído à flexibilidade podia ser, na
verdade, benefício de uma política de reserva mais frouxa.

A diferença medida no dia de vento alto é de 472 para 300\,MW de carga de data center vista pelo
requisito, ou seja, 172\,MW a menos, que a 3\,\% dão cerca de 5\,MW de reserva a menos, contra um $R_t$
médio de 378\,MW. É pouco, 1{,}4\,\%, mas o argumento de um revisor não seria sobre a magnitude e sim
sobre a assimetria existir.

Implementei então a opção \texttt{reserve.common\_reference}, que soma de volta o equivalente plano da
energia flexível em C e D, igualando $R_t$ nos quatro casos, e reresolvi C e D nos quatro dias.

\begin{table}[htbp]
\centering\small
\caption{Grade base contra requisito de reserva idêntico nos quatro casos (k\$/dia).}
\tabreserva
\end{table}

O efeito existe, mas é pequeno. O valor do data center cai de \numdCplano{} para
\numdCresref{}\,k\$/dia, ou 1{,}6\,\%, bem dentro do ruído do gap, que é de 2{,}1 a 2{,}3\,k\$. O
benefício conjunto cai de 99{,}6 para \numdDresref{}, e \textbf{a interação fica praticamente igual}, de
\numIbase{} para \numIresref{}\,k\$/dia, sem mudar de sinal em nenhum dos quatro dias.

Onde o efeito aparece de verdade é nas partidas, que é justamente a métrica em que o data center mais
contribui. A redução de partidas do uso conjunto cai de \numredsubase{} para \numredsuDresref{}\,\%. Ou
seja, \textbf{cerca de 3 pontos percentuais dos \numredsubase{}\,\% vinham do requisito de reserva mais
folgado}, e não da flexibilidade em si. É pouco, uns 7\,\% do que se atribuía ao recurso, mas é honesto
reportar.

A conclusão do trabalho não muda, e agora existe o número para responder a quem perguntar. Mantive a
grade base como principal, porque ela representa a comparação operacionalmente correta, na qual carga
adiável de fato não precisa ser coberta por reserva girante, e esta entra como sensibilidade.

\subsection{Carga e descarga simultâneas}

Na formulação livre, nenhuma restrição impede que uma bateria carregue e descarregue na mesma hora; em
tese, ela poderia dissipar energia nas perdas para reduzir curtailment. A verificação \emph{a posteriori}
separa dois casos. Nas baterias da VPL, que pagam degradação de 2\,\$/MWh, a sobreposição é exatamente
nula em todas as soluções: dissipar nunca compensa. Já a unidade de armazenamento nativa do RTS-GMLC
(50\,MW no barramento 313), que não tem custo de degradação no \emph{deck}, opera com sobreposição em 3
a 6 horas por dia, dissipando de 18 a 23\,MWh/dia nas perdas para absorver excedente renovável. O efeito
é praticamente idêntico nos quatro casos (151\,MWh/dia em A e B, 141 em C e 119 em D), então se cancela
nas diferenças: responde por 32 dos 2\,407\,MWh de redução de curtailment atribuída ao uso conjunto, ou
1{,}3\,\%, e por cerca de 2\,\$/dia no custo. Não altera nenhuma conclusão, mas é uma degenerescência do
sistema teste que vale registrar.

% =====================================================================
\section{Conclusões}

\subsection{Resultado principal}

O enunciado mais preciso do resultado não é ``VPL e VDC são substitutos'', mas:

\textbf{Dois recursos de flexibilidade capazes de aliviar o mesmo gargalo competem economicamente mesmo
quando entregam serviços operacionais bem diferentes} --- um reduz sobretudo curtailment, o outro sobretudo
ciclagem térmica.
O enunciado é sustentado pelos dados e tem consequência prática direta para dimensionamento.

\subsection{Implicações práticas}

\begin{enumerate}
\item \textbf{Para o planejamento, que é o terreno da tese}: dimensionar armazenamento de corredor e carga
flexível de forma independente superestima o benefício conjunto --- de 1 a 11\,\% no porte base, e até um
quarto do benefício da carga flexível quando a bateria é grande. A métrica $I$, que sai de quatro rodadas
day-ahead, é um número que um estudo de dimensionamento poderia incorporar.
\item \textbf{Para a conexão de data centers}: ligá-los como carga firme joga fora um recurso de redução
de ciclagem térmica que, nas minhas simulações, equivale ao de uma bateria de 200\,MW. Ressalva: as tarefas são sintéticas.
\item \textbf{Para quem modela}: reportar o ruído do gap sempre que medir interação entre recursos.
\end{enumerate}

\subsection{Resultados contrários à expectativa}

\begin{itemize}
\item O pareamento vale pouco na operação: quase todo o benefício vem da localização das baterias nos
terminais corretos, que é justamente o que a tese resolve.
\item O data center atua principalmente sobre a ciclagem, mas sua economia \emph{monetária} vem três quartos de
combustível e só um quarto de custo de partida, porque as unidades que ele desliga são turbinas a gás
baratas. O benefício de ciclagem é real, mas subvalorizado pelo objetivo econômico convencional.
\item O spread de preço não colapsa, se desloca.
\end{itemize}

% =====================================================================
\section{Limitações}

\begin{itemize}
\item Determinístico e sem segurança N-1; a incerteza aparece só por análise de sensibilidade.
\item Quatro dias representativos mais doze mensais, não um ano completo.
\item As tarefas são tratadas como divisíveis e interrompíveis, o que é um limite superior da
flexibilidade real.
\item Os data centers são sintéticos em tamanho, PUE e perfil de carga.
\item Economia \emph{operacional}, sem custo de investimento. Responde ``quanto vale operar'', não ``vale
a pena instalar''.
\end{itemize}

% =====================================================================
\section{Pontos para discussão}

O modelo está implementado em PyPSA, com cerca de 3\,700 linhas de código próprio, tudo reproduzível por
um \texttt{Makefile}, e com cinco testes automatizados. Os resultados deste relatório saem todos do mesmo
arquivo de métricas, sem nenhum número digitado à mão.

Pontos para discussão:

\begin{enumerate}
\item \textbf{O enquadramento.} Dado que o pareamento explica só 4\,\% do benefício, faz sentido continuar
falando em ``Virtual Power Line'', ou seria mais honesto falar em ``armazenamento nos terminais do
corredor''? A opção que considero mais adequada é manter VPL, pela conexão com a linha do grupo e porque a regra da
tese é testada, mas gostaria da sua avaliação.
\item \textbf{A regra de chaveamento custar de 2 a 2{,}4\,\%.} O resultado pode ser lido como crítica ao
trabalho do Flávio, o que não é a intenção. Mantenho o enquadramento como custo de oportunidade de um
protocolo simples em modelo idealizado, ou o senhor sugere outro tratamento?
\item \textbf{O dia de referência.} As tabelas e figuras usam o dia de vento alto, onde o mecanismo do corredor é mais visível. O dia mediano tem apenas uma hora de saturação e não evidenciaria o mecanismo. Concorda com a escolha?
\item \textbf{O próximo passo.} A extensão natural é o UC estocástico, para verificar se a
competição se mantém sob incerteza. A alternativa é aplicar ao sistema brasileiro com dados reais. Qual
das duas é mais útil para a tese?
\end{enumerate}

\end{document}
